\chapter{Introduction}

\section{Context and Motivation}

Retrieval-Augmented Generation (RAG) has become a default way to give
large language models access to information beyond their initial
training data. Instead of relying solely on parametric knowledge, a
RAG-based model retrieves relevant data from an external corpus at
inference time and adjusts its responses based on that retrieved
context. This pattern is now used in a variety of production systems
across search, customer support, and enterprise knowledge tools:
enterprise assistants that answer questions against internal
documentation, customer-support chatbots that ground replies in a
company's own knowledge base, and consumer-facing assistants that cite
sources for factual claims rather than relying purely on what a model
memorized during training. The appeal is straightforward. Parametric
knowledge freezes at a model's training cutoff and cannot reflect
anything that happened afterward, and a model reasoning purely from its
own weights has no built-in mechanism to point to a specific source for
a claim it makes. Grounding generation in retrieved text addresses both
problems at once, at the cost of introducing a new component into the
pipeline that an adversary can target: the retrieval corpus itself.

Retrieval has also opened a new attack pattern. An adversary who can
influence the retrieved content or the conversations built around it can
manipulate the model's output without ever touching its model weights,
without needing any access to the model's training process, and often
without the system's operators noticing anything unusual about the
request itself. The European Union's Artificial Intelligence Act
(Regulation (EU) 2024/1689) has given this question a regulatory
importance: Article 15 requires high-risk AI systems to display "an
appropriate level of accuracy, robustness, and cybersecurity" and to be
resilient against attempts to exploit system vulnerabilities. This
dissertation treats Article 15(5) as a technical specification to test
against, not a legal claim to argue, a distinction maintained throughout
the rest of the dissertation.

\section{Problem Statement}

Three distinct mechanisms exist for manipulating a RAG system's output
through its retrieval layer, and treating them as variations on one idea
rather than three separate threats is itself a methodological risk this
dissertation avoids. The first, planting a malicious instruction inside
retrieved text, known as indirect prompt injection \cite{greshake2023not},
exploits the fact that a model has no reliable way to distinguish an
instruction written by its system prompt from one embedded in a document
it retrieves. The second, planting a false fact into the model's corpus
itself, known as knowledge poisoning \cite{zou2025poisonedrag}, targets
the corpus rather than the model's instruction-following behavior: an
adversary inserts a small number of crafted documents into the knowledge
base such that, for a chosen query, the poisoned documents are retrieved
and steer the model toward a specific false answer. The third,
manipulating the model across multiple conversational turns, exploiting
its conversational memory state rather than any single retrieved
document \cite{russinovich2025great}, escalates a benign-seeming
conversation gradually toward an output the model would refuse if asked
directly in a single turn. Each targets a genuine vulnerability in the
pipeline rather than a variation on a single idea, so a defense effective
against one offers no guarantee against another.

Existing RAG security benchmarks tend to cover a broader set of attacks
within a single attack category or one narrow slice of the threat
surface, rather than all three families side by side under one shared,
controlled environment. SafeRAG \cite{liang2025saferag} evaluates four
RAG-specific attack tasks (noise, conflict, toxicity, denial-of-service)
across fourteen RAG components, but its safety framework has no
multi-turn conversational attacks, and it does not test or evaluate the
same four target models this research studies. RSB \cite{zhang2025rsb}
benchmarks thirteen poisoning attacks and seven defenses at far greater
scale than this research's single poisoning method, but covers poisoning
exclusively, with no injection or multi-turn attack family at all. RSB
remains an arXiv preprint at the time of writing this, and no other
trusted version has emerged despite deep searches across the web. This
dissertation includes it here as an explicit exception to its
peer-review sourcing standard because it was one of three benchmark
works confirmed directly with the supervisor. Crucially, its own primary
finding shows that its strongest defense reduces poisoning success but
at an accuracy drop exceeding 20\% in most settings, which parallels
this dissertation's own finding that spotlighting is the only defense
with a confirmed real effect on injection, at a mean $\Delta F_1$ of
$-0.283$. Hidden-in-Plain-Text \cite{guo2026hidden} reports
attack-success confidence intervals for indirect prompt injection
specifically, but is scoped to web-carrier delivery mechanisms and
doesn't cover knowledge poisoning or multi-turn jailbreaking. No
existing work to this dissertation's knowledge evaluates all three
attack families against the same set of models under a single
backend-controlled harness; the backend-confound discovery detailed in
Chapter 4 is itself a methodological consideration most of this
literature doesn't appear to control for.

\section{Research Questions}

This dissertation explores two research questions:

\begin{itemize}
    \item \textbf{RQ1:} How vulnerable are small-to-mid-size (4--12B
    parameters) open-source, open-weight, instruction-tuned language
    models to indirect prompt injection, knowledge poisoning and
    multi-turn adversarial manipulation when deployed in a
    retrieval-augmented generation pipeline?
    \item \textbf{RQ2:} How effective are practical, low-cost defenses
    at reducing attack success and at what cost to the system's
    underlying answer utility?
\end{itemize}

A third, narrower question follows from the first two and runs through
the closing chapters of this dissertation rather than standing on its
own: how do the answers to RQ1 and RQ2 map onto Article 15(5)'s specific
vulnerability classes, and where does that mapping expose gaps in what a
benchmark of this kind can and cannot demonstrate? This third question
is framed deliberately as an open mapping exercise rather than a claim
to be proven, since the goal of Chapter 5's regulatory discussion is
honesty about where the benchmark's evidence is strong and where it is
silent, not forcing every result into a compliance verdict it cannot
support.

\section{Contributions}

This dissertation contributes:

\begin{enumerate}
    \item A reproducible, statistically rigorous, real-hardware
    benchmark spanning four open-source language models, two retrieval
    corpora, three attack families and three defenses, built so that
    every reported number traces back to a specific, re-runnable script
    and a verified result file rather than a one-off measurement.
    \item A documented account of an inference-backend confound
    discovered mid-analysis, which overturned two of the three defenses'
    apparent headline effects once corrected, offering a citable
    methodological warning for anyone building a similar evaluation: a
    defense's apparent effect can be, in whole or in part, an artifact
    of which inference backend its baseline happened to run on, and this
    is detectable only by deliberately building a backend-matched
    comparison.
    \item A mapping of these results onto Article 15(5)'s vulnerability
    classes, offered as a concrete technical contribution to a
    compliance-testing gap that currently has no harmonised standard,
    together with an explicit account of which parts of Article 15 this
    benchmark's design cannot speak to.
\end{enumerate}

The central empirical finding running through all three contributions is
that robustness is not one factor. Model identity dominates injection
success. Corpus identity dominates knowledge-poisoning success. The
multi-turn attack produces no statistically significant ranking across
models at all. A benchmark, defense, or compliance argument built around
one blended notion of robustness would miss this structure entirely, and
would risk recommending a defense calibrated to the wrong axis of
vulnerability.

\section{Implementation and Reproducibility}

This dissertation evaluates four open-source, instruction-tuned language
models: Llama-3.1-8B-Instruct \cite{grattafiori2024llama3}, Qwen3-8B
\cite{yang2025qwen3}, Phi-4-mini-instruct \cite{abouelenin2025phi4mini},
and Ministral-3-8B-Instruct-2512 \cite{mistralai2025ministral3}. These
four models were selected against five explicit requirements, detailed
further in Chapter 4: a permissive open license, self-hostable weights,
instruction tuning, a parameter count in the roughly 4--12B range, and a
dense, non-mixture-of-experts architecture. Every result rests on real
generations run on real GPU hardware \cite{runpod2024platform}, not
simulated or estimated output, a distinction that matters given how much
of the adjacent literature relies on smaller samples or proxy
measurements. The full implementation, an open-source Python framework,
is available at
\url{https://github.com/Simarjit1303/rag-adversarial-robustness}
\cite{singh2026ragrepo}; it satisfies this module's technical-contribution
requirement and lets another researcher reproduce, extend, or challenge
any result reported here directly against the same code and data.

\section{Scope and Limitations}

This dissertation uses two retrieval corpora, HotpotQA
\cite{yang2018hotpotqa} and MS MARCO \cite{nguyen2016msmarco}, across
every comparison. It also evaluated a third corpus, Natural Questions
Open \cite{kwiatkowski2019natural}, and excluded it after it became
structurally unsuitable for testing retrieval-dependent robustness: its
passage text is, by construction, the gold answer itself, a property
detailed with full evidence in Chapter 5. The supervisor confirmed this
decision as a settled scope choice. This dissertation excludes model
poisoning and confidentiality attacks from its scope: every attack it
studies operates at inference time, and none touch training data or
model weights, a boundary worth stating plainly rather than stretching
any of the three studied attacks to cover ground they were never
designed for. It does not determine whether any specific deployed RAG
system counts as high-risk under Annex III of the AI Act; that
determination depends on deployment context outside this benchmark's
scope, and this dissertation makes no claim about it in either
direction.

\section{Structure of the Dissertation}

This chapter has presented the motivation, problem, research questions,
and contributions towards this research. Figure~\ref{fig:dissertation-structure}
outlines the chapters that follow.

\begin{figure}[htbp]
\centering
\resizebox{\linewidth}{!}{%
\begin{tikzpicture}[
  box/.style={rectangle, draw, rounded corners, align=center, minimum width=2.1cm, minimum height=1.3cm, font=\small},
  arr/.style={-Stealth, thick}
]
\node[box] (c1) {Ch.~1\\Introduction};
\node[box, right=0.5cm of c1] (c2) {Ch.~2\\Foundations};
\node[box, right=0.5cm of c2] (c3) {Ch.~3\\Related Work};
\node[box, right=0.5cm of c3] (c4) {Ch.~4\\Approach};
\node[box, right=0.5cm of c4] (c5) {Ch.~5\\Evaluation};
\node[box, right=0.5cm of c5] (c6) {Ch.~6\\Conclusion};
\draw[arr] (c1) -- (c2);
\draw[arr] (c2) -- (c3);
\draw[arr] (c3) -- (c4);
\draw[arr] (c4) -- (c5);
\draw[arr] (c5) -- (c6);
\end{tikzpicture}%
}
\caption{Structure of the dissertation, chapters in reading order.}
\label{fig:dissertation-structure}
\end{figure}

As Figure~\ref{fig:dissertation-structure} shows, five chapters follow
after this one. Chapter 2 (Foundations) establishes the technical
background needed to follow the rest of the dissertation: retrieval-
augmented generation itself, the model-versus-context distinction that
organizes this dissertation's whole attack taxonomy, and the regulatory
context this dissertation tests against. Chapter 3 (Related Work)
surveys the existing literature on RAG-specific adversarial attacks and
defenses, organized by attack family and defense mechanism; under the
framework set by the university's CDS structure, this chapter serves as
the dissertation's literature review. Chapter 4 (Approach) describes the
benchmark's design and implementation in detail, including the
infrastructure it runs on and the real engineering failures encountered
while building it. Chapter 5 (Evaluation and Results) presents the full
statistical analysis across all three attack families and three
defenses, and maps them onto Article 15. Chapter 6 (Conclusion)
summarizes the findings, states their limitations honestly, and closes
with a discussion of what this benchmark contributes to, and cannot
resolve about, the compliance-testing question this dissertation opened
with.